%-------------------------------------------------
% FileName: chap-2.tex
% Description: 第2章 相关技术与理论基础
%-------------------------------------------------

\chapter{相关技术与理论基础}

\section{前端开发技术}

\subsection{Vue3}

本系统前端使用 Vue3。选择它并不是因为它是主流，而是因为项目页面之间存在着比较明显的状态差别，用户端有登录、场地列表、预约、器材租借、个人订单等页面，管理端有场地维护、订单管理、用户管理、数据看板。直接用 HTML 和零星的脚本来实现页面切换、表单回显、分页刷新等逻辑，会使页面的组织变得凌乱。Vue3 的组件化方式可以把每一个页面拆分成独立的视图，接口请求则集中在 api 目录下，代码结构也会更加清晰\cite{wang2020}。

Vue3 的响应式机制主要是用在几类操作上。预约页面根据场地开放时间产生可选时段，用户切换日期之后再次请求 /api/venues/\{id\}/schedule，页面会自动更新已约时段和可选结束时间；订单页在支付或者取消之后重新加载列表，状态标签随着返回数据的变化；后台看板使用 computed 统计开放场地数、已支付订单数、器材可借率。这些地方若用手工操作 DOM，则代码会变得繁杂，Vue3 可以将数据的变化同界面的更新绑定在一起，适合于本系统这类交互并不复杂但是状态较多的业务页面。

\subsection{Element Plus}

Element Plus 是给 Vue3 用的一套较常见的组件库。本项目中能看见它的地方主要是表格、表单、弹窗、分页、消息提示这五类——后台页面全部使用了 el-table 和 el-dialog，前台的下单、租借也是用 el-form 配合校验规则来实现数据的收集；操作反馈全部通过 ElMessage 和 ElMessageBox 来完成\cite{zhang2019}。这些控件的样式以及交互行为较为成熟，可以省去自己写组件、调样式的时间，对于一个学生级的项目来说特别合适。

风格上的好处是统一性。整个系统从前台到后台使用相同的配色、控件，按钮、表格、提示弹层都不会有前后页面不同造成的违和感，上手更简单。

\section{后端开发技术}

\subsection{Spring Boot}

后端采用 Spring Boot 2.7。它之所以被选中是因为 starter 依赖已经将 web、validation、jdbc 等常用东西集成了起来，消除了 Spring 早期版本大段 XML 配置的需求，项目按 controller、service、mapper 三层分了开来，每个模块的登录、场地、订单、器材等模块都在一个包内，新人接下来看代码可以快速找到对应部分\cite{luo2019}。

在实现上，Spring Boot 只做“基建”，即接口暴露、参数绑定、参数校验和全局异常处理等工作。看起来很简单的问题，却需要在每个控制器里重新写相同的参数判空和统一返回包装代码，整体可读性就会下降很多。MyBatis-Plus、JWT 工具类、数据库连接池等组件使用 Spring 的依赖注入方式一起工作，实际开发业务的时候只需要关注自己的需求就可以了。

\subsection{JWT 认证机制与接口设计}

登录鉴权选用了 JWT。它的特点就是无状态，服务端不需要保存会话，token 本身就已经包含有用户 ID 和角色信息，后端拦截器拿到 token 解析之后就可以直接判断了。这种结构在前后端分离的模式下比较方便，也不必在多实例环境下还要进行 session 共享\cite{wu2018}。本系统的做法是登录成功后把 token 通过响应头返回前端，后续请求由 Axios 拦截器统一塞进 Authorization 头，再由后端 JwtInterceptor 解析。

接口设计上遵循 RESTful 风格，路径按认证、场地、订单和管理端等业务域划分，所有响应都用一个统一的 Result<T> 对象包装，里面是 code、message、data 三个字段。这样的好处是前端处理成功和失败可以走同一条分支：先看 code，再决定怎么用 data。

\section{数据存储与开发环境}

\subsection{MySQL 数据库技术}

数据存储用的是 MySQL 8.0。对于一个校园规模的系统来说，关系型数据库可以满足基本的要求，社区资料也比较齐全，调试、排错都很方便。本系统的数据模型是以角色、用户、场地、订单、器材、租借记录六个表为基础构建的，外键以及状态字段把这些实体联系在一起，从而形成完整的业务语义\cite{zhang2017}。

数据库不仅仅负责信息的存储，预约冲突判断、订单状态转换、库存扣减等逻辑也被直接放在了 SQL 里实现，冲突判断用的是按照 venue\_id 和 book\_date 进行范围查询再加唯一索引，库存扣减则放进了 @Transactional 中防止并发导致的数量出现负数。这些决定着业务能否正常运转。

\subsection{系统开发与运行环境}

开发环境的具体版本是 Java 17、Spring Boot 2.7、MyBatis-Plus 3.5、MySQL 8.0、Vue3 + Vite、Node.js 18。后端使用 Maven 管理依赖，前端使用 Vite 启动开发服务，并进行热更新以及构建。组合是目前 Web 项目中比较常见的配置，社区活跃，遇到问题很容易找到现成的方案。

部署时前后端可以分别独立运行，前端构建产物是一组静态文件，用 Nginx 或任意静态托管即可，后端打包成 jar 文件直接 java -jar 启动，数据库独立部署。这样一种结构使得三者可以分阶段地添加或者缩减，当压力过大时可以单独增加一台后端的机器而不需要动用全部系统的资源。
